\begin{helper}[Boundary selection outside the base tree]
Let $X_0=V(T)$ and $X=\noderef{ScheideEnlargedVertexSet}(R)$.  If $X$ is
not strongly closed, then there are a color $\alpha$, a defective-color
index $d=(\delta,h_\delta)$, an edge $f$, and vertices $w,u$ such that
\[
 \alpha\ne\alpha_0,\qquad \varphi(f)=\alpha,\qquad
 G.\mathop{\rm ends}(f)=s(w,u),
\]
\[
 w\in X\setminus X_0,\qquad u\notin X,\qquad
 w=(R d).u_1\ \hbox{or}\ w=(R d).u_2.
\]
Moreover, for the full $(\alpha_0,\alpha)$ incidence component
\[
 Q=\noderef{TwoColorComponent}(G,\varphi,\alpha_0,\alpha,
      \operatorname{inl}u),
\]
the edge-node $\operatorname{inr}f$ belongs to $Q$, and no vertex-node
$\operatorname{inl}v$ with $v\in X_0$ belongs to $Q$.
\end{helper}
\begin{proof}
Write $y_0=C.y(0)$.  The global maximum triple contained in
\noderef{ScheideStructuralConfiguration}, via
\noderef{GlobalMaximumTriple} and \noderef{KTriple}, says that $\varphi$
is proper.  The base-position equation puts $y_0$ in
$X_0=\noderef{TreeSequenceVertexSet}(T)$, and
\noderef{ScheideDefectiveSetup} says that $\alpha_0$ is missing at
$y_0$.  By \noderef{ScheideGlobalMaximumOrder}, $X_0$ is closed and has
odd order, and by \noderef{ScheideEnlargedClosed}, $X$ is closed.

Since $X$ is not strongly closed, the definition
\noderef{ColoredVertexSetStronglyClosed} gives a color $\alpha$ on two
distinct identities of \noderef{BoundaryEdges}$(G,X)$.  Closedness of
$X$ implies that $\alpha$ is present at every vertex of $X$: otherwise a
vertex missing $\alpha$ and either boundary identity contradict
\noderef{ColoredVertexSetClosed}.  In particular, the missingness of
$\alpha_0$ at $y_0\in X$ and the definition \noderef{PresentColors}
give $\alpha\ne\alpha_0$.

At each vertex of $X$ there is exactly one incident $\alpha$-edge.
Existence is presentness.  Two distinct such identities would be
\noderef{AdjacentEdges} and violate properness.  An $\alpha$-edge in
\noderef{InternalEdges}$(G,X)$ contributes two incidences, an
$\alpha$-edge in the boundary contributes one, and all other edges
contribute none; this is immediate from \noderef{Incident} and the
distinct-end law in \noderef{LooplessMultigraph}.  Hence
\[
 |X|=2|E_\alpha(X)|+|\partial_{\varphi,\alpha}(X)|.       \tag{1}
\]

For every defective $\delta$, the internal edge supplied by
\noderef{ScheideSymmetricDefectiveInternal} has the two distinct vertices
$(R\delta).u_1,(R\delta).u_2$ as its ends.  Both lie outside $X_0$ by
\noderef{ScheideReplacementData}.  The equal-or-disjoint alternative in
\noderef{ScheideReplacementPairsLaminar} therefore says that the indexed
union in \noderef{ScheideEnlargedVertexSet} adds a disjoint union of
two-vertex sets to $X_0$.  Thus $|X|$ is odd.  Equation (1) makes the
number of $\alpha$-boundary identities odd; since it is at least two, it
is at least three.

At most one of these identities has its end in $X_0$.  If $\alpha$ is
not in \noderef{DefectiveColors}$(G,\varphi,X_0)$, this is the definition
of nondefectiveness.  Suppose instead that $\alpha$ is defective and use
it as an index $a$.  The first clause of
\noderef{ScheideSymmetricDefectiveInternal} gives an internal
$\alpha$-edge on the deleted pair $U_{\gamma(\alpha)}$.  Properness
excludes every other $\alpha$-edge at either deleted vertex, so each
$\alpha$-boundary edge of $X_0$ leaves from the retained set
$Z_{\gamma(\alpha)}$.  The exact-three clause of
\noderef{ScheideReplacementData} identifies those identities as
$f^0_\alpha,f^1_\alpha,f^2_\alpha$.  The latter two end in the two
distinct vertices of $U_\alpha$, hence are internal to
$X_\alpha=X_0\cup U_\alpha$.  Properness at these vertices also shows
that the outside end of $f^0_\alpha$ is not in $U_\alpha$.  Consequently
$f^0_\alpha$ is the sole $\alpha$-boundary identity of $X_\alpha$, so it
is the only one among the three that can leave $X\supseteq X_\alpha$.
Thus at least two $\alpha$-boundary identities of $X$ have their in-end
in $X\setminus X_0$.

Let $P$ be the full $(\alpha_0,\alpha)$ component through $y_0$.
Because $\alpha$ is present throughout $X$,
\noderef{ScheideSingleBoundaryCrossing} gives exactly one edge-node of
$P$ in $\partial_G(X)$.  It has color $\alpha$: every edge-node in $P$
has one of the two component colors by \noderef{TwoColorComponent} and
\noderef{TwoColorIncidenceGraph}, while closedness of $X$ and the missing
color $\alpha_0$ at $y_0$ forbid an $\alpha_0$ boundary edge.  Select,
among the two boundary identities whose in-end is outside $X_0$, an edge
$f$ different from this unique identity.  Orient its distinct ends as
$G.\mathop{\rm ends}(f)=s(w,u)$ with
$w\in X\setminus X_0$ and $u\notin X$; this orientation follows from
\noderef{BoundaryEdges} and looplessness.  Expanding
\noderef{ScheideEnlargedVertexSet} supplies a defective index $d$ for
which $w=(R d).u_1$ or $w=(R d).u_2$.

Let $Q$ be the full $(\alpha_0,\alpha)$ component through $u$.  Since
$f$ has color $\alpha$ and is incident with $u$, the incidence edge
definition puts $\operatorname{inr}f$ in $Q$.  The components $P$ and
$Q$ are distinct, because otherwise $f$ would be the boundary edge-node
of $P$; connected components are therefore disjoint.

It remains to prove that $Q$ avoids $X_0$.  No $\alpha_0$-edge leaves
$X_0$, by its closedness and the missingness of $\alpha_0$ at $y_0$.  If
$\alpha$ is nondefective, its at most one boundary identity of $X_0$ lies
on $P$: the component begins at $y_0\in X_0$ and reaches the unique
boundary crossing of $X$, hence its first exit from $X_0$ must use that
identity.  A first entry of the disjoint component $Q$ into $X_0$ would
use either a forbidden $\alpha_0$ boundary identity or that same
$\alpha$ identity, a contradiction.

If $\alpha$ is defective, then
\noderef{ScheideReplacementUnionClosed} says that $X_\alpha$ is closed.
It contains $y_0$, so it has no $\alpha_0$ boundary identity, and the
argument above showed that its unique $\alpha$ boundary identity is
$f^0_\alpha$.  The component $P$ starts in $X_\alpha$ and reaches outside
$X\supseteq X_\alpha$, so its first exit is $f^0_\alpha$.  Since $Q$
starts outside $X$ and is disjoint from $P$, the same first-entry argument
prevents $Q$ from entering $X_\alpha$.  In either case no vertex-node of
$X_0$ belongs to $Q$, as required.
\end{proof}
